\chapter{Social dance}
% mini introduction about the chapter
\begin{chapterabstract}
	Následující kapitolu uvedu do problematiky společenských tanců, vysvětlím co to jsou společenské tance, jaké druhy společenských tanců existují a jaké jsou jejich charakteristiky. Následně zanalyzuji jaké nástroje a technologie jsou využívány uživateli pro vyhledávání kurzů a akcích spojenými s těmito tanci. Nakonec popíši motivaci pro vytvoření aplikace, která by uživatelům to vyhledávání usnadnila. 
\end{chapterabstract}
\section{Co jsou to společenské tance}
Společenské tance (více známé jako nepřekládané "social dance") jsou tance, které se tančí především kvůli zábavě a společenské interakci mezi lidmi, na rozdíl od jiných forem tanců, které se zaměřují na soutěžení. Avšak je důležité poznamenat, že některé 
taneční akce nebo festivaly začali zahrnovat i možnost zúčastnit se soutěží. Soutěže ovšem bývají jinak postavené než klasické taneční soutěže, protože partneři bývají náhodně přiřazeni. Tudíž tanečníci nemají možnost si s partnerem předem nacvičit choreografii, ale musí se spolehnout na své schopnosti improvizace a komunikace během tance.
\cite{socialDance:about}

Společenské tance bývají nejčastěji párové tance, kde partneři spolu komunikují prostřednictvím pohybů a vedení. Jeden z partnerů obvykle bývá označován jako "vedoucí" (leader) a druhý jako "následující" (follower). Leader vede tanec a určuje směr a styl pohybu, zatímco follower reaguje na jeho vedení a doplňuje tanec svými vlastními pohyby. Co je na tanci zajímavé, je to že tyto tance se tancují většinou improvizovaně, což znamená, že neexistují pevně stanovené kroky nebo choreografie, které by museli tanečníci dodržovat. Místo toho leader vymýšlí kroky a pohyby na místě a pomocí správného vedení umožňuje followerovi správně reagovat. 
 \cite{socialDance:about}

\section{Druhy společenských tanců}
Existuje mnoho druhů společenských tanců, které se liší svým stylem, rytmem a kulturou, ze které pocházejí. Já jsem se rozhodla pro účely mé diplomové práce se zaměřit na druhy tanců se kterými se nejvíce setkávám v mé taneční komunitě. Jedná se o následující tance:
\begin{itemize}
    \item Salsa
    \item Bachata
    \item Zouk
\end{itemize}
V následující sekci popíši charakteristiky těchto tanců. Samozřejmě můžeme zmínit i další tance, se kterými se člověk taky může setkat, jako například Kizomba, Swing, Reggaeton a další, ale pro účely této práce jsem se rozhodla zaměřit pouze na ty tři, které jsem zmínila výše, protože jsou mi nejvíce blízské a setkávám se s nimi nejčastěji.

Také jsem se setkala s tím, že se taneční styly kombinují dohromady, tomuto pojmu se říká "fusion". \cite{centrumtance_dance_fusion}

Nejvíce jsem se za poslední dobu setkala s trendem, kde se kombinuje bachata s zoukem, tento styl se nazývá "bachazouk". TODO

\subsection{Salsa}
Je latinsko-americký tanec s kubánskými kořeny vyznačující se rychlým tempem, energickými pohyby a výraznou rytmikou. Tanec působí živě a radostně, což z něj činí vhodnou volbu pro společenské události a taneční večery. \cite{socialDance:salsa}

Tento tanec vznikl zhruba ve 30 letech 20. století na Kubě a později se rozšířil do dalších částí Latinské Ameriky. Salsa se následně kvůli migraci rozšířila i do Spojených států, kde byla braná jako populární zábava. Díky tomuto se salsa různě proměňovala a vyvíjela, což vedlo k vytvoření různých stylů salsy. \cite{socialDance:salsa:history}

Salsa se většinou tančí na hudbu ve čtyřčtvrtečním taktu na doby 1, 2, 3 a 5, 6, 7, přičemž doba 4 a 8 jsou pauzy (avšak některé styly salsy mohou mít odlišné rytmus jako například New York salsa, která se tančí na doby 2, 3, 4 a 6, 7, 8). \cite{socialDance:salsa:rhythm}

Z mé zkušenosti je salsa často na akcích ruku v ruce s bachatou, například v jedné místnosti se tančí salsa a v druhé místnosti se tančí bachata. Také se dějí tzv. SBK party, což znamená, že se tančí salsa, bachata a kizomba. 

\subsection{Bachata}
Bachat je tanec pocházející z Dominikánské republiky, který je rozšířený po celém světě. 
Je známá především jako sensuální tanec. V tomto tanci se tančí sekvence kroků na 8 dob, v originálu se tančí do čtverce, ale v západním světě se začalo do strany, což je jednodušší. V obou případech se tančí 3 kroky a poté následuje tzv. "tap", což je 
krok, kdy se noha zvedne a poklepe o zem (často doprovázená bokem do strany), ale nedělá se krok. \cite{incognitodance_what_is_bachata}

Bachata se dělí na několik stylů, přičemž tanečníci často kombinují různé styly dohromady. Chtěla bych zmínit ty nejběžnější styly, které se tančí v mé komunitě, jako například \cite{jtdancestudio_bachata_styles}:
\begin{itemize}
    \item Dominikánská bachata - tento styl je ve své původní podobě jednou z nejjednodušších, ale zároveň nejživějších forem tohoto tance. Tento styl se vyznačuje hravým footworkem a pohybem boků do strany podle rytmu hudby. Tento styl neobsahuje komplikované figury, což ho činí vynikajícím pro začátečníky.
    \item Bachata Moderna - nazývána také jako "moderní bachata" je styl, který mixuje tradiční bachatu s jinými tanečními styly, jako je například salsa. Tento styl se především vyznačuje většími pohyby, dynamičtější projev a rychlé otočky. Také obsahuje více záklonů než bychom třeba shledali u původní bachaty.
    \item Sensual - tento styl vznikl ve Španělsku díky tanečníkům Korke Escalona a Judith Cordero za cílem přidat do bachaty více výrazů a emocí. Tento styl kombinuje prvky moderní bachaty s kroky brazilského zouku. Typickými prvky toho stylu jsou vlnící se pohyby těla a izolace (oddělené pohyby různých částí těla). Tanečníci obvykle tančí blízko sebe a pohybují se jako jeden celek, což navozuje emotivní a romantickou atmosféru.
\end{itemize}

\subsection{Zouk}
Je partnerský tanec, pocházející z Brazílie. Tento tanec vznikl v počátcích 90. let 20. století a byl ovlivněn karabskou zouk hudbou a tancem lambada. Během let se zouk vyvinul do několika stylů, přičemž se dá tančit nejenom na zouk hudbu, ale i například na hip hop, R\&B, pop hudbu a další. Zouk se obecně vyznačuje vedením pomocí celého těla, nakloněné (říká se off-axis) pozice, otočky a náklony partnerky. Charakteristické pro zouk jsou také pohyby hlavy partnerky (říká se jim tzv. “hair movement”), které vznikají nakláněním followera při vedení. Tyto pohyby jsou dle mého názoru jedním z nejzajímavějších aspektů tohoto tance a zároveň i jedním z nejtěžších, protože vyžadují velkou kontrolu nad tělem a správné vedení. Zároveň jsem si všimla, že tyto figury se začaly objevovat i v bachtě, díky tomu, jak vypadají dobře. \cite{zoukminneapolis_brazilian_zouk}



\section{Cesta tanečníka ve světě společenských tanců}
Lidé se do světa společenských tanců často dostávají z různých důvodů. Z mojí zkušenosti nejčastěji prostřednictvím přátel, kteří již tančí, sociálních sítí, kde jsou propagovány taneční workshopy a kurzy, nebo prostřednictvím videí na platformách jako je YouTube, kde lidé objevují různé taneční styly a chtějí se je naučit.

Většinou si člověk vybere konkrétní tanec, který ho zaujal, a začne hledat kurzy nebo taneční školy, které tento tanec nabízejí. Tudíž prvním krokem bývá nalezení vhodného kurzu.
Po pár lekcích tanečníci často začínají navštěvovat taneční akce, kde mohou tančit s různými partnery a zlepšovat své dovednosti v reálném prostředí. Na kurz člověk nemusí přijít v páru, protože se na lekcích partneři většinou střídají. Občas jsem se i setkala s tím, že kurz bývá koncipován tak, že nejprve proběhne výuková hodina/workshop, a poté následuje 
party, při které lidé procvičují co se naučili a improvizují. Zkušenější tanečníci se mohou účastnit různých tanečních festivalů, v rámci nichž probíhají workshopy a taneční večery. Některé taneční školy nebo organizace nabízejí i sérii workshopů anebo hodin v rámci kurzu, na které se mohou tanečníci jednotlivě hlásit a tudíž si mohou vybrat konkrétní lekce, které je zajímají anebo se jim hodí do časové kapacity, místo toho, aby se hlásili na celý kurz.
Velkým zážitkem dle mého názoru je i cestování na taneční akce do jiných měst nebo zemí, kde mohou tanečníci poznat nové lidi a kultury prostřednictvím tance. \cite{socialDance:about}

\section{Existující prostředky pro vyhledávání událostí}
V předchozí sekci jsem popsala jak se tanečníci dostávají do světa společenských tanců a jaké události mohou navštěvovat. V této sekci zanalyzuji jaké nástroje a technologie jsou využívány uživateli pro vyhledávání kurzů a akcích spojenými s těmito tanci. Při tvorbě této analýzy jsem vycházela z vlastních zkušeností s vyhledáváním kurzů a akcí, zejména v oblasti bachaty v Praze, jelikož dlouhodobě na tyto akce docházím a jsem schopna lépe zhodnotit dostupné nástroje. Tato analýza byla zároveň konzultována s dalšími tanečníky působícími v této komunitě, kteří poskytli podněty vedoucí k jejímu zlepšení.

\subsection{Webové stránky tanečních škol}
Webové stránky byly pro mě prvním místem, kde jsem hledala informace o kurzech a akcích. Většina tanečních škol má své vlastní webové stránky, kde zveřejňují informace o nabízených kurzech, termínech, cenách a dalších detailech. Tyto stránky často obsahují i kalendář akcí, kde jsou uvedeny nadcházející taneční večery a workshopy. Nicméně, nevýhodou je, že uživatel
musí procházet webové stránky tanečních škol jednotlivě a stranou si poznamenávat informace o kurzech, což může být časově náročné a únavné. Navíc je možné, že nějaké organizace/spolky nemají vlastní webové stránky, ale působí jen na sociálních sítích, tudíž na ně uživatel nemůže narazit na webu.

Existují i stránky, které agregují informace z více tanečních škol a spolků. Některé tyto stránky jsou koncipované tak, že abyste mohli přidat událost, musíte uzavřít spolupráci s danou stránkou, což často zahrnuje i poplatek anebo si berou provizi z prodeje vstupenek. Toto může být pro některé organizace odrazující faktor, což vede k tomu, že událostí na těchto stránkách není tak hojný počet. Jedna z příkladů těchto stránek je go \& dance, která se hlavně zaměřuje na propagaci velkolepých festivalů a kongresů. Tyto akce bývají v řádech stovek eur, což je poměrně drahé v porovnáním kolik stojí běžné kurzy a taneční večery. Stránka má hezký design a je přehledná, avšak pokud se snažím najít akce konající se v Praze, tak mi stránka vrátí i akce konající se v jiných městech. Toto je pro mě poměrně frustrující, protože musím procházet všechny akce a hledat ty, které se konají v Praze.

Jako další příklad jsem nalezla stránku tanci.cz, která cílila na kurzy a události pořádané v ČR, avšak tato stránka se jeví jako nefunkční a nenalézá mi žádné kurzy.
%% mam uvadet obrazkama priklady?
\subsection{Sociální sítě}
Sociální sítě jsou dle mého pozorování nejběžnějším prostředkem k vyhledávání nadcházejících tanečních akcích a zároveň i slouží jako platforma pro komunikaci mezi tanečníky a organizátory. Dle mého pozorování je nejčastější platformou pro tento účel Facebook. 

Na facebooku organizátoři vytvářejí události, kde zveřejňují informace o nadcházejících akcích, jako jsou termíny, místa konání, ceny a další detaily. Většina z nich pro registrace používá buď že napíšete zprávu organizátorovi anebo google formuláře. Registrace bývá občas trochu chaotická, protože vyplníte například google formulář, ale nemáte jistotu, že organizátor váš zájem o účast zaznamenal, protože vám nepřijde žádné potvrzení o registraci. Také se stává, že tím jak je organizování dáno přes zprávy, tak se populárním organizátorům špatně udržují čekací listiny. To může vést k demotivaci lidí, kteří se snaží do kurzu dostat.  

Uživatelé mohou události na facebooku sledovat, přihlásit se na ně a sdílet je se svými přáteli. Facebook také umožňuje organizátorům cíleně propagovat své události prostřednictvím placených reklam, což zvyšuje jejich viditelnost a dosah. Nicméně, nevýhodou je, že pokud sociální síť nevyhodnotí, že by daná akce byla pro uživatele relevantní, nezobrazí ji. Tím uživatel může přijít o informace o akcích, které by ho mohly zajímat.  

Uživatelé musí procházet různé skupiny, stránky a události, aby našli relevantní informace o nadcházejících akcích, což může být časově náročné a frustrující. Je také docela těžké dohledat akci co v případě hledání třeba klíčového slova "bachata" umístěného v Praze, v názvu bachata nemá. Pak mi facebook akci nenabídne, ale nabízí mi akce, které mají v názvu "bachata", ale nejsou umístěné v Praze.
Z mé zkušenosti velmi pomohlo, když jsem se na facebooku propojila s tanečníky, protože mi to umožnilo vidět události, o které mají zájem a tímto se stalo pro mne jednodušší akce dohledávat. 

Další sociální sítí, která se využívá pro tento účel, je Instagram. Na Instagramu organizátoři často sdílejí informace o nadcházejících akcích prostřednictvím příspěvků a příběhů. Uživatelé mohou sledovat účty organizátorů a získávat informace o nadcházejících akcích. Nicméně je zde ještě horší vyhledávání akcí a spíše se instagram využívá pro propagaci akcí a komunikaci s uživateli, než pro samotné vyhledávání.

Po té jsou tu i další platformy jako například Whatsapp, kde se často vytvářejí skupiny pro tanečníky, kteří spolu komunikují a sdílejí informace o nadcházejících akcích. Tyto skupiny jsou často uzavřené a přístupné pouze pro členy, což může být nevýhodou pro nové tanečníky, kteří se snaží najít informace o akcích.

Tyto sociální sítě jsou velmi užitečné pro komunikaci a propagaci akcí, ale nejsou ideální pro samotné vyhledávání, protože sociální sítě jsou zaměřené obecně na komunikaci a sdílení, a ne na organizaci a vyhledávání tanečních akcí. další problémem se kterým se potýkám je rozptýlenost informací.
Některé akce jsou propagovány pouze na facebooku, jiné pouze na instagramu a další pouze prostřednictvím whatsappových skupin. Tudíž dochází k tomu, že uživatel musí procházet různé platformy, aby našel informaci, kterou někdy viděl, ale už si nepamatuje kde. Dále tím, že jedna akce je na více platformách, tak se může stát, že jsou informace o ní nekonzistentní. Například se mi jednou stalo, že jsem se přihlásila na akci na facebooku a později jsem se přidala do whatsappové skupiny, která souvisela s touto akcí, ale účast v ní nebyla povinná pro účast na akci. Po té došlo ke změně času konání akce, ale já jsem se o této změně dozvěděla až v den konání akce, ale informace byla zprostředkována pouze prostřednictvím whatsappové skupiny, ale ne na facebooku. Tím že většina lidí o akci věděla pouze z facebooku, tak se stalo, že mnoho lidí se na akci dostavilo pozdě a nemohli se tak zúčastnit naplánovaného workshopu. dalším problémem je, že uživatel kromě svých přátel nevidí kdo se na akci přihlásil a jaký je poměr mezi leadery a followery.

\subsection{Mobilní aplikace}
Mobilních aplikací, které by zahrnovali různé taneční akce, je poměrně málo. Při prohlížení jednotlivých aplikací jsem narážela na problém, že neobsahují skoro žádné akce a jeví se jako poměrně neaktivní. Poté mi jeden známý doporučil aplikaci "Only Dance", která se nejenom zaměřuje na taneční akce ale také funguje jako sociální síť pro tanečníky. Tato aplikace má poměrně hezký design a nabízí i funkcionality jako možnost vytvářet spojení mezi tanečníky nebo psát příspěvky, které celá komunita může vidět. Nicméně, nevýhodou je, že aplikace obsahuje pouze možnost přidat reakci na akci, že plánuju jít anebo že se mi akce líbí, ale neobsahuje možnost rozlišit role pro akce, kde by to bylo relevantní. Také nemohu vidět kdo na akci jde, pokud uživatele nemám ve spojení. Dále mi přijde škoda, že uživatel nemůže přidávat videa na profil, určitě by to usnadnilo hledání tanečního partnera. Co mě na aplikaci nepříjemně překvapilo je, že když se akce zruší, tak ji aplikace smaže ale neoznačí ji jako zrušenou, což může být pro uživatele matoucí, protože s akcí již počítal a najednou se mu ztratí z přehledu, aniž by věděl, že se akce zrušila.

\subsection{Shrnutí - motivace pro vytvoření nové aplikace}
V předchozích sekcích jsem popsala různé nástroje a technologie, které jsou využívány pro vyhledávání kurzů a akcí spojenými se společenskými tanci. Z mého pohledu je zřejmé, že existuje několik problémů, které uživatelé mohou zažít při hledání těchto informací. Tyto problémy zahrnují rozptýlenost informací, nekonzistentní informace o akcích, nedostatek přehledu o tom, kdo se na akci přihlásil a zda uživatelova akce vedla k úspěšné registraci.